\documentclass{article}
\usepackage{booktabs}
\usepackage{geometry}
\geometry{margin=1in}

\title{\textbf{Problem Set 10:} 
\\Cross-Validation and Classification}
\author{}
\date{April 21, 2026}

\begin{document}
\maketitle

\section*{Overview}
This problem set compares five machine learning algorithms on the UCI Adult Income dataset. The goal is to classify whether an individual earns more than \$50,000 per year (\texttt{high.earner}). Each algorithm was tuned via 3-fold cross-validation, and final performance was evaluated on a held-out test set comprising 20\% of the data.

\section*{Optimal Tuning Parameters and Out-of-Sample Performance}

Table \ref{tab:results} reports the optimal hyperparameter values selected by cross-validation and the resulting out-of-sample accuracy for each algorithm.

\begin{table}[h!]
\centering
\caption{Optimal Tuning Parameters and Test Set Accuracy}
\label{tab:results}
\begin{tabular}{llll}
\toprule
Algorithm & Tuning Parameter(s) & Optimal Value(s) & Test Accuracy \\
\midrule
Logistic Regression & penalty ($\lambda$) & 0.0000000001 & 0.853 \\
Decision Tree & cost\_complexity & 0.001 & 0.868 \\
                    & tree\_depth & 15 & \\
                    & min\_n & 10 & \\
Neural Network & hidden\_units & 7 & 0.845 \\
               & penalty ($\lambda$) & 1 & \\
k-Nearest Neighbor & neighbors & 30 & 0.843 \\
SVM (RBF) & cost & 1 & 0.864 \\
          & rbf\_sigma & 0.25 & \\
\bottomrule
\end{tabular}
\end{table}

\section*{Comparison of Out-of-Sample Performance}

The decision tree achieved the highest out-of-sample accuracy at 86.8\%, followed closely by the SVM at 86.4\%. Logistic regression performed well at 85.3\%, despite being the simplest model. The neural network and k-nearest neighbor classifier performed slightly worse, at 84.5\% and 84.3\% respectively.

Several observations are worth noting. First, the decision tree's optimal cost complexity was very low (0.001), suggesting the data supports a relatively complex tree structure. The optimal depth of 15 with a minimum node size of 10 indicates the tree was allowed to grow fairly deep. Second, the SVM selected a moderate cost ($C=1$) and a small kernel width ($\sigma=0.25$), reflecting a relatively smooth decision boundary. Third, the logistic LASSO selected a near-zero penalty, meaning regularization was minimal and most predictors were retained. Fourth, the neural network favored 7 hidden units with strong regularization (penalty = 1), suggesting some complexity but with shrinkage to prevent overfitting. Finally, kNN selected the maximum number of neighbors tested (30), indicating that smoother, more averaged decision boundaries worked best.

Overall, the tree-based and kernel-based methods outperformed the distance-based and linear approaches on this dataset, though the differences in accuracy are modest (less than 3 percentage points across all five algorithms).

\end{document}